% REV02: pruebas reunidas en 02_dualidad_t.tex y 03_pantallas_coordinadas.tex.
\subsection{Incidence exceptionnelle et coordination des écrans}
\label{exc:pantallas-dualidad}

La direction sélectionnée par \(K\) et le réseau de Leech proviennent
de représentations distinctes de l’incidence construite. La première
détermine la réduction orthogonale \(12\to11\to10\) ; le second
intervient dans la réduction isotrope \(26\to24\). La coordination
conserve l’espace préalable, ses relations et l’information transmise
par chaque projection.

Les sections sur la dualité T et les écrans coordonnés de \cite{hmtX} développent conjointement cette coordination. La première réunit l'involution des coordonnées entières, la conjugaison unitaire avec ses domaines et le changement de section résiduelle avec mémoire du quotient. La seconde construit le morphisme \(\mathcal R_M^{\mathrm{alg}}\) qui y est défini et démontre l'isomorphisme des présentations quotientées dans le théorème correspondant sur les écrans isomorphes. Le graphe conserve la coordonnée dans \(V_{11}\) avec sa projection dans \(V_{10}\), tandis que le quotient isotrope récupère le réseau de Leech. Les noyaux et codomaines respectifs sont ainsi préservés.

L’ellipse d’action est construite avant d’appliquer la dualité et fixe
son rayon adimensionnel ainsi que le rayon réciproque
(section de l’ellipse d’action de \cite{hmtX}). Le secteur compact relie alors
l’incidence exceptionnelle à une forme quadratique déterminée. Les
champs, l’action et les opérateurs de la réalisation physique ultérieure
sont présentés dans leur domaine propre, en conservant cet antécédent
algébrique et ses démonstrations internes.
